% 2.总体分析.tex —— 二、总体分析
\section{总体分析}
本题的本质是求解无限长圆柱内的\textbf{径向一维非稳态热质传递}问题。热量由烘房经表面对流换热带入药材内部，水分由内部向表面迁移并经表面汽化进入烘房。二者在几何上共用同一圆柱体，在时间上共用同一环境条件记录，但在题设附录 2 的参数体系下，热方程不含水分项（未给出汽化潜热），故热与质之间为\textbf{单向耦合}，两方程可按“先热后质”的顺序逐时间步求解。

数值方法的选择是本题的核心。本文采用“先显式、后隐式”的递进策略：首先\textbf{完整实现 Euler 显式方法}，把径向节点分为中心点、内部点、边界点三类分别离散，其优点是逐点推进、无需解线性方程组、对非线性 $D(C)$ 无需迭代，缺点是受 CFL 稳定条件限制、时间只有一阶精度；在此基础上\textbf{引入 Crank--Nicolson 隐式格式与 Thomas 追赶法}，用无条件稳定性换取时间二阶精度与更大的时间步。这一递进既能清楚呈现差分格式的构造逻辑，又能说明隐式格式相对显式格式的定量收益。
